% ============================================================
% Margin Requirements — Simplified Rule-Based Framework
% Trade Republic — Margin Trading Methodology
% Compile with: pdflatex margin_methodology_simple.tex
% ============================================================

\documentclass[9pt, a4paper]{extarticle}

% ── Packages ─────────────────────────────────────────────
\usepackage[a4paper, top=2.5cm, bottom=2.5cm, left=3cm, right=3cm]{geometry}
\usepackage{amsmath, amssymb}
\usepackage{booktabs}
\usepackage{array}
\usepackage{tabularx}
\usepackage{titlesec}
\usepackage{enumitem}
\usepackage{fancyhdr}
\usepackage{hyperref}
\usepackage{microtype}
\usepackage[T1]{fontenc}
\usepackage[sfdefault]{inter}
\usepackage{parskip}
\usepackage{longtable}
\usepackage{textcomp}
\usepackage{xcolor}
\usepackage{pdflscape}
\usepackage{colortbl}

% ── Hyperref ─────────────────────────────────────────────
\hypersetup{
    hidelinks,
    pdftitle  = {Margin Requirements — Simplified Rule-Based Framework},
    pdfauthor = {},
}

% ── Section heading styles ────────────────────────────────
\titleformat{\section}
    {\large\bfseries}
    {\thesection.}{0.6em}{}
    [\titlerule]

\titleformat{\subsection}
    {\normalsize\bfseries}
    {\thesubsection}{0.6em}{}

\titleformat{\subsubsection}
    {\normalsize\itshape}
    {\thesubsubsection}{0.6em}{}

\titlespacing{\section}{0pt}{1.6em}{0.8em}
\titlespacing{\subsection}{0pt}{1.2em}{0.4em}
\titlespacing{\subsubsection}{0pt}{1.0em}{0.3em}

% ── Header / Footer ──────────────────────────────────────
\pagestyle{fancy}
\fancyhf{}
\renewcommand{\headrulewidth}{0.4pt}
\fancyhead[L]{\small\textit{Margin Requirements — Simplified Rule-Based Framework}}
\fancyhead[R]{\small\textsc{Strictly Private \& Confidential} $|$ April 2026}
\fancyfoot[C]{\small\thepage}

% ── Formula block environment ─────────────────────────────
\newenvironment{formulabox}{\par\addvspace{4pt}}{\par\addvspace{4pt}}

% ── Variable table command ────────────────────────────────
\newcommand{\vartableheader}{%
    \textbf{Symbol} & \textbf{Description}\\
}

% ── Euro symbol ──────────────────────────────────────────
\newcommand{\EUR}[1]{\ifmmode\text{\texteuro}#1\else\texteuro#1\fi}

% ── Alert box ────────────────────────────────────────────
\newcommand{\alertbox}[1]{%
    \par\vspace{4pt}%
    \noindent\colorbox{gray!12}{\parbox{\dimexpr\linewidth-2\fboxsep}{%
        \small #1}}%
    \vspace{4pt}\par%
}

% ── Misc ─────────────────────────────────────────────────
\setlength{\parindent}{0pt}
\setlength{\parskip}{5pt}

% ─────────────────────────────────────────────────────────
\begin{document}

% ── Title page ───────────────────────────────────────────
\begin{titlepage}
\vspace*{4cm}
\begin{center}
    {\Huge\bfseries Margin Requirements}\\[0.5em]
    {\Large Simplified Rule-Based Framework}\\[2em]
    {\large Trade Republic Bank GmbH}\\[0.5em]
    {\large Risk Modelling}\\[2em]
    {\normalsize April 2026}\\[1em]
    {\small\textsc{Strictly Private \& Confidential}}
\end{center}
\end{titlepage}

\tableofcontents
\newpage

% ─────────────────────────────────────────────────────────
\section{Overview}
\label{sec:overview}
% ─────────────────────────────────────────────────────────

This document specifies a simplified, rule-based margin framework for margin trading
products. Margin requirements are expressed as fixed percentage rates per asset class,
applied directly as minimum equity conditions on the outstanding loan. No econometric
modelling is required.

The framework rests on a single unifying principle: the \textbf{haircut is the margin
requirement}. An initial margin haircut $h_{\mathrm{IM}}$ of, say, 30\% means the
customer must hold equity equal to at least 30\% of the current market value at all
times during the IM phase. Equivalently, the outstanding loan must not exceed
$(1 - h_{\mathrm{IM}})$ times the market value. This identity connects the haircut
directly to the LTV monitoring limit.

The framework consists of four margin components:

\begin{enumerate}[leftmargin=2em]
    \item \textbf{Initial Margin (IM)} — minimum equity condition at position open,
          enforced for one trading day (the IM horizon). Expressed as rate $h_{\mathrm{IM}}$
          on market value.
    \item \textbf{Maintenance Margin (MM)} — ongoing minimum equity condition from
          day two onwards. Expressed as rate $h_{\mathrm{MM}} \leq h_{\mathrm{IM}}$ on
          market value, giving the customer more headroom once the initial risk window
          has passed.
    \item \textbf{Overnight Margin} — additional equity cushion required at end-of-day
          (EOD) to cover the overnight risk window. Expressed as add-on rate
          $h_{\mathrm{ON}}$ applied on top of the prevailing phase rate.
    \item \textbf{Mark-to-Market (MTM) Margin} — continuous revaluation of the
          outstanding loan against current market prices. Updates the LTV in real time;
          a breach of the applicable LTV limit triggers immediate auto-liquidation.
\end{enumerate}

% ─────────────────────────────────────────────────────────
\section{Core Definitions}
\label{sec:definitions}
% ─────────────────────────────────────────────────────────

\subsection{Market Value and Loan}

Let $N$ denote the number of units held in instrument~$i$ and $P_t$ the current
market price. The \textbf{market value} at time~$t$ is:
\[
    \mathrm{MV}_t = N \times P_t
\]

The \textbf{outstanding loan} $L_t$ is the principal borrowed plus accrued interest
to date. Customer equity is:
\[
    E_t = \mathrm{MV}_t - L_t
\]

\subsection{Loan-to-Value Ratio}

The \textbf{LTV ratio} measures the loan relative to the current market value:

\begin{formulabox}
\[
    \mathrm{LTV}_t = \frac{L_t}{\mathrm{MV}_t}
\]
\end{formulabox}

An LTV of zero means the position is fully equity-funded. An LTV of one means the
customer has no equity remaining. Because margin trading limits leverage, the LTV is
always strictly below one.

\subsection{The Haircut as LTV Limit}

The core identity of this framework is that \textbf{the haircut $h$ defines both the
margin requirement and the LTV monitoring limit}:

\begin{formulabox}
\[
    \text{Margin requirement:} \quad E_t \;\geq\; h \times \mathrm{MV}_t
\]
\[
    \text{Equivalent LTV condition:} \quad \mathrm{LTV}_t \;=\; \frac{L_t}{\mathrm{MV}_t}
    \;\leq\; 1 - h
\]
\end{formulabox}

The LTV limit for a given phase is therefore $1 - h$, where $h$ is the applicable
margin rate. A higher margin rate produces a lower (more conservative) LTV limit.

\begin{tabularx}{\linewidth}{lX}
\toprule
\textbf{Symbol} & \textbf{Description} \\
\midrule
$h_{\mathrm{IM}}$   & Initial margin rate; applies on day of opening ($t = 0$) \\
$h_{\mathrm{MM}}$   & Maintenance margin rate; applies from $t = 1$ onwards \\
$h_{\mathrm{ON}}$   & Overnight margin add-on; applied at EOD on top of prevailing phase rate \\
$T_{\mathrm{IM}}$   & IM horizon: 1 trading day \\
$L_t$               & Outstanding loan at time $t$ \\
$\mathrm{MV}_t$     & Market value at time $t$ \\
$\mathrm{LTV}_t$    & Loan-to-Value ratio at time $t$ \\
\bottomrule
\end{tabularx}

% ─────────────────────────────────────────────────────────
\section{Margin Components}
\label{sec:components}
% ─────────────────────────────────────────────────────────

\subsection{Initial Margin (IM)}

Initial Margin applies at the moment a position is opened and remains in force for
one trading day ($T_{\mathrm{IM}} = 1$). It is the most conservative condition,
ensuring the customer enters the position with sufficient equity to absorb an adverse
intraday move.

\begin{formulabox}
\[
    \text{IM condition:} \quad \mathrm{LTV}_0 = \frac{L_0}{\mathrm{MV}_0} \;\leq\; 1 - h_{\mathrm{IM}}
\]
\end{formulabox}

The maximum loan at opening and the minimum required equity follow directly:
\[
    L_{\max} = (1 - h_{\mathrm{IM}}) \times \mathrm{MV}_0
    \qquad\qquad
    E_{\min} = h_{\mathrm{IM}} \times \mathrm{MV}_0
\]

The auto-liquidation price during the IM phase (the price at which
$\mathrm{LTV} = 1 - h_{\mathrm{IM}}$ exactly) is:
\[
    P^{*}_{\mathrm{IM}} = \frac{L_0}{N \times (1 - h_{\mathrm{IM}})}
\]

\subsection{Maintenance Margin (MM)}

After one trading day, the position transitions from the IM phase to the Maintenance
Margin phase. The MM rate $h_{\mathrm{MM}} \leq h_{\mathrm{IM}}$ is less conservative,
raising the LTV limit and giving the customer more room. The same structural condition
applies:

\begin{formulabox}
\[
    \text{MM condition:} \quad \mathrm{LTV}_t = \frac{L_t}{\mathrm{MV}_t} \;\leq\; 1 - h_{\mathrm{MM}}
    \qquad (t \geq 1)
\]
\end{formulabox}

The auto-liquidation price in the MM phase:
\[
    P^{*}_{\mathrm{MM}} = \frac{L_t}{N \times (1 - h_{\mathrm{MM}})}
\]

Because $h_{\mathrm{MM}} < h_{\mathrm{IM}}$, we have $1 - h_{\mathrm{MM}} > 1 -
h_{\mathrm{IM}}$ and therefore $P^{*}_{\mathrm{MM}} < P^{*}_{\mathrm{IM}}$: the
position can absorb a larger adverse price move before liquidation fires in the MM
phase.

\subsection{Overnight Margin}

At end-of-day, an additional overnight add-on $h_{\mathrm{ON}}$ is applied on top of
the prevailing phase rate. This reflects the extended risk window during which
continuous intraday monitoring is unavailable.

\begin{formulabox}
\[
    \text{Overnight condition:} \quad \mathrm{LTV}^{\mathrm{ON}}_t = \frac{L_t}{\mathrm{MV}_t}
    \;\leq\; 1 - h_{\mathrm{phase}} - h_{\mathrm{ON}}
\]
\end{formulabox}

where $h_{\mathrm{phase}}$ is $h_{\mathrm{IM}}$ on day~0 and $h_{\mathrm{MM}}$ from
day~1 onwards. The overnight LTV limit is therefore tighter (lower) than the intraday
limit by exactly $h_{\mathrm{ON}}$.

If the overnight condition is not met at EOD, the customer must either:
\begin{enumerate}[leftmargin=2em]
    \item deposit additional cash to reduce the LTV below $1 - h_{\mathrm{phase}} - h_{\mathrm{ON}}$, or
    \item partially close the position before the EOD cut-off.
\end{enumerate}
Failure to remedy the breach triggers automatic liquidation of the excess exposure.

The overnight auto-liquidation price:
\[
    P^{*}_{\mathrm{ON}} = \frac{L_t}{N \times (1 - h_{\mathrm{phase}} - h_{\mathrm{ON}})}
\]

\subsection{Mark-to-Market (MTM) Margin}

MTM reflects the continuous revaluation of the open position as market prices move.
It is not a separate cash margin call but a real-time update to the LTV:

\begin{formulabox}
\[
    \mathrm{MTM}_t = \bigl(P_t - P_{t-1}\bigr) \times N
\]
\[
    \mathrm{LTV}_t = \frac{L_t}{N \times P_t}
\]
\end{formulabox}

A negative MTM (adverse price move) raises the LTV; a positive MTM lowers it. The
LTV is recomputed \textbf{continuously} at each price update throughout the trading
day. As soon as $\mathrm{LTV}_t$ breaches the applicable limit, auto-liquidation is
triggered immediately — there is no grace period.

\alertbox{\textbf{Auto-liquidation rule:} at any time $t$, if $\mathrm{LTV}_t > 1 - h_{\mathrm{phase}}$
(intraday) or $\mathrm{LTV}_t > 1 - h_{\mathrm{phase}} - h_{\mathrm{ON}}$ (overnight check),
the position is automatically liquidated at market prices. No margin call is issued.}

% ─────────────────────────────────────────────────────────
\section{Two-Phase Monitoring Framework}
\label{sec:monitoring}
% ─────────────────────────────────────────────────────────

\subsection{Phase Transition}

The IM horizon is exactly one trading day. The transition from IM to MM is automatic:

\[
    h_{\mathrm{phase}}(t) =
    \begin{cases}
        h_{\mathrm{IM}} & t = 0 \text{ (day of opening)} \\
        h_{\mathrm{MM}} & t \geq 1
    \end{cases}
\]

\subsection{Intraday Monitoring (Both Phases)}

At every price update the LTV is recomputed and checked against the intraday limit:

\[
    \mathrm{LTV}_t = \frac{L_t}{N \times P_t} \;\leq\; 1 - h_{\mathrm{phase}}(t)
\]

\subsection{EOD Overnight Check (Both Phases)}

At end-of-day, the overnight add-on tightens the limit:
\[
    \mathrm{LTV}^{\mathrm{ON}}_t = \frac{L_t}{N \times P_t^{\mathrm{close}}}
    \;\leq\; 1 - h_{\mathrm{phase}}(t) - h_{\mathrm{ON}}
\]

\subsection{LTV Limits Summary}

\begin{tabularx}{\linewidth}{lllX}
\toprule
\textbf{Phase} & \textbf{Time} & \textbf{LTV limit} & \textbf{Action on breach} \\
\midrule
IM — intraday   & $t = 0$     & $1 - h_{\mathrm{IM}}$                          & Auto-liquidation \\
IM — overnight  & $t = 0$ EOD & $1 - h_{\mathrm{IM}} - h_{\mathrm{ON}}$        & Top-up / reduce or auto-liquidation \\
MM — intraday   & $t \geq 1$  & $1 - h_{\mathrm{MM}}$                          & Auto-liquidation \\
MM — overnight  & $t \geq 1$ EOD & $1 - h_{\mathrm{MM}} - h_{\mathrm{ON}}$     & Top-up / reduce or auto-liquidation \\
\bottomrule
\end{tabularx}

% ─────────────────────────────────────────────────────────
\section{Asset Class Rate Table}
\label{sec:rates}
% ─────────────────────────────────────────────────────────

Fixed margin rates are maintained per asset class. All rates are expressed as a
percentage of current market value.

\begin{tabularx}{\linewidth}{lXXX}
\toprule
\textbf{Asset Class}
  & \textbf{$h_{\mathrm{IM}}$}
  & \textbf{$h_{\mathrm{MM}}$}
  & \textbf{$h_{\mathrm{ON}}$} \\
\midrule
Major index equities (e.g.\ EuroStoxx~50, S\&P~500 constituents) & 15\% & 10\% & 3\% \\
Other equities                                                     & 30\% & 25\% & 5\% \\
Equity ETFs — broad market                                         & 15\% & 10\% & 3\% \\
Equity ETFs — sector / thematic                                    & 20\% & 15\% & 4\% \\
Investment-grade bonds ($\geq$ BBB)                                & 10\% & 7\%  & 2\% \\
High-yield bonds ($<$ BBB)                                         & 30\% & 20\% & 5\% \\
\bottomrule
\end{tabularx}

\bigskip
\textit{Rates are reviewed at least quarterly. Individual instruments may be assigned
higher rates at the discretion of Risk if recent volatility or liquidity conditions
warrant it.}

% ─────────────────────────────────────────────────────────
\section{Worked Example}
\label{sec:example}
% ─────────────────────────────────────────────────────────

\subsection{Setup}

A customer opens a leveraged long position in an ``Other equity'' (non-index stock).
Applicable rates: $h_{\mathrm{IM}} = 30\%$, $h_{\mathrm{MM}} = 25\%$,
$h_{\mathrm{ON}} = 5\%$.

\begin{tabularx}{\linewidth}{lX}
\toprule
\textbf{Parameter} & \textbf{Value} \\
\midrule
Current price $P_0$         & \EUR{100.00} per share \\
Number of shares $N$        & 150 \\
Market value $\mathrm{MV}_0$ & $150 \times \EUR{100} = \EUR{15{,}000}$ \\
Customer equity $E_0$       & \EUR{5{,}000} \\
Outstanding loan $L_0$      & \EUR{15{,}000} $-$ \EUR{5{,}000} $= \EUR{10{,}000}$ \\
Leverage                    & $3{\times}$ \\
\bottomrule
\end{tabularx}

\bigskip

Applicable LTV limits:

\begin{tabularx}{\linewidth}{llll}
\toprule
\textbf{Phase} & \textbf{Rate(s)} & \textbf{LTV limit} & \textbf{Liquidation price $P^{*}$} \\
\midrule
IM intraday    & $h_{\mathrm{IM}} = 30\%$                        & $1 - 0.30 = 70\%$                           & $\EUR{10{,}000}/(150 \times 0.70) = \EUR{95.24}$ \\
IM overnight   & $h_{\mathrm{IM}} + h_{\mathrm{ON}} = 35\%$     & $1 - 0.35 = 65\%$                           & $\EUR{10{,}000}/(150 \times 0.65) = \EUR{102.56}$ \\
MM intraday    & $h_{\mathrm{MM}} = 25\%$                        & $1 - 0.25 = 75\%$                           & $\EUR{10{,}000}/(150 \times 0.75) = \EUR{88.89}$ \\
MM overnight   & $h_{\mathrm{MM}} + h_{\mathrm{ON}} = 30\%$     & $1 - 0.30 = 70\%$                           & $\EUR{10{,}000}/(150 \times 0.70) = \EUR{95.24}$ \\
\bottomrule
\end{tabularx}

\subsection{Step 1 — Opening Check (IM Phase, $t = 0$)}

\[
    \mathrm{LTV}_0 = \frac{\EUR{10{,}000}}{\EUR{15{,}000}} = 66.7\%
    \;\leq\; 70\%
    \quad\Rightarrow\quad \textbf{Position may be opened.}
\]

Minimum equity the customer must post:
\[
    E_{\min} = h_{\mathrm{IM}} \times \mathrm{MV}_0 = 30\% \times \EUR{15{,}000} = \EUR{4{,}500}
\]
Customer posted \EUR{5{,}000} $>$ \EUR{4{,}500}: condition satisfied with a buffer of \EUR{500}.

The IM intraday auto-liquidation price:
\[
    P^{*}_{\mathrm{IM}} = \frac{L_0}{N \times (1 - h_{\mathrm{IM}})}
    = \frac{\EUR{10{,}000}}{150 \times 0.70}
    = \mathbf{\EUR{95.24}}
\]
The stock must fall $-4.76\%$ from \EUR{100} before intraday auto-liquidation triggers.

\subsection{Step 2 — EOD Overnight Check ($t = 0$ EOD)}

The overnight LTV limit tightens to $1 - h_{\mathrm{IM}} - h_{\mathrm{ON}} = 65\%$.
At EOD price \EUR{100} (unchanged intraday):
\[
    \mathrm{LTV}^{\mathrm{ON}}_0 = \frac{\EUR{10{,}000}}{\EUR{15{,}000}} = 66.7\%
    \;>\; 65\%
    \quad\Rightarrow\quad \textbf{Overnight condition breached.}
\]

The overnight auto-liquidation price:
\[
    P^{*}_{\mathrm{ON}} = \frac{\EUR{10{,}000}}{150 \times 0.65} = \EUR{102.56}
\]
At the current price of \EUR{100} the LTV already exceeds the 65\% overnight limit.
The customer must post additional collateral before EOD cut-off.

\textbf{Required top-up} to clear the overnight condition exactly:
\[
    L_{\max}^{\mathrm{ON}} = 65\% \times \EUR{15{,}000} = \EUR{9{,}750}
    \qquad\Rightarrow\qquad
    \text{Cash to post} = L_0 - L_{\max}^{\mathrm{ON}} = \EUR{10{,}000} - \EUR{9{,}750} = \EUR{250}
\]
After posting \EUR{250}: $L = \EUR{9{,}750}$, LTV$^{\mathrm{ON}} = 65\% = $ limit. Passed.

\subsection{Step 3 — Transition to MM Phase ($t = 1$)}

After one trading day the IM phase ends and MM governs. The LTV limit widens from
70\% to 75\% (intraday). Assuming price remains \EUR{100} and loan is \EUR{9{,}750}
after the overnight top-up:

\[
    \mathrm{LTV}_1 = \frac{\EUR{9{,}750}}{\EUR{15{,}000}} = 65.0\%
    \;\leq\; 75\%
    \quad\Rightarrow\quad \textbf{MM phase check passed.}
\]

The new MM intraday auto-liquidation price:
\[
    P^{*}_{\mathrm{MM}} = \frac{\EUR{9{,}750}}{150 \times 0.75} = \mathbf{\EUR{86.67}}
\]
The position now tolerates a $-13.3\%$ price decline before auto-liquidation — significantly
more room than the IM phase.

\subsection{Step 4 — Intraday MTM: Price Falls $-5\%$ to \EUR{95}}

MTM is calculated continuously. At price \EUR{95}:
\[
    \mathrm{MTM} = (95 - 100) \times 150 = -\EUR{750}
    \qquad
    \mathrm{MV} = 150 \times \EUR{95} = \EUR{14{,}250}
\]
\[
    \mathrm{LTV} = \frac{\EUR{9{,}750}}{\EUR{14{,}250}} = 68.4\%
    \;\leq\; 75\%
    \quad\Rightarrow\quad \textbf{Still open.}
\]

\subsection{Step 5 — Intraday MTM: Price Falls Further to \EUR{88}}

\[
    \mathrm{MV} = 150 \times \EUR{88} = \EUR{13{,}200}
\]
\[
    \mathrm{LTV} = \frac{\EUR{9{,}750}}{\EUR{13{,}200}} = 73.9\%
    \;\leq\; 75\%
    \quad\Rightarrow\quad \textbf{Still open.}
\]

\subsection{Step 6 — Intraday MTM: Price Falls to \EUR{86.50}}

\[
    \mathrm{MV} = 150 \times \EUR{86.50} = \EUR{12{,}975}
\]
\[
    \mathrm{LTV} = \frac{\EUR{9{,}750}}{\EUR{12{,}975}} = 75.1\%
    \;>\; 75\%
    \quad\Rightarrow\quad \textbf{AUTO-LIQUIDATION TRIGGERED.}
\]

The position is automatically liquidated at market price \EUR{86.50}. The customer
receives the residual equity:
\[
    E_{\mathrm{final}} = \EUR{12{,}975} - \EUR{9{,}750} = \EUR{3{,}225}
\]
The customer has lost $\EUR{5{,}000} - \EUR{3{,}225} = \EUR{1{,}775}$ (35.5\% of their
initial equity) due to the $-13.5\%$ price decline.

\subsection{Summary Table}

\begin{tabularx}{\linewidth}{lllllX}
\toprule
\textbf{Event} & \textbf{Price} & \textbf{Loan} & \textbf{MV} & \textbf{LTV} & \textbf{Limit / Status} \\
\midrule
Open ($t=0$, IM intraday)      & \EUR{100}    & \EUR{10{,}000} & \EUR{15{,}000} & $66.7\%$ & $\leq 70\%$ \quad \textbf{Open} \\
EOD overnight ($t=0$)          & \EUR{100}    & \EUR{10{,}000} & \EUR{15{,}000} & $66.7\%$ & $\leq 65\%$ \quad \textbf{Breach} \\
After top-up \EUR{250}         & \EUR{100}    & \EUR{9{,}750}  & \EUR{15{,}000} & $65.0\%$ & $\leq 65\%$ \quad \textbf{Cleared} \\
MM phase open ($t=1$)          & \EUR{100}    & \EUR{9{,}750}  & \EUR{15{,}000} & $65.0\%$ & $\leq 75\%$ \quad \textbf{Open} \\
MTM $-5\%$ ($t=1$)             & \EUR{95}     & \EUR{9{,}750}  & \EUR{14{,}250} & $68.4\%$ & $\leq 75\%$ \quad \textbf{Open} \\
MTM $-12\%$ ($t=1$)            & \EUR{88}     & \EUR{9{,}750}  & \EUR{13{,}200} & $73.9\%$ & $\leq 75\%$ \quad \textbf{Open} \\
MTM $-13.5\%$ ($t=1$)         & \EUR{86.50}  & \EUR{9{,}750}  & \EUR{12{,}975} & $75.1\%$ & $> 75\%$ \quad \textbf{LIQUIDATE} \\
\bottomrule
\end{tabularx}

% ─────────────────────────────────────────────────────────
\newpage
\begin{landscape}

\section{Comparison with Alternative Margin Frameworks}
\label{sec:comparison}

The table below applies all four frameworks to the same position:
\textbf{150 shares at \EUR{100}, market value \EUR{15{,}000}, customer equity
\EUR{5{,}000}, outstanding loan \EUR{10{,}000} ($3{\times}$ leverage)}.
Where a framework does not permit the full $3{\times}$ position, the maximum
feasible position on \EUR{5{,}000} equity is shown instead.

\definecolor{highlight}{RGB}{255,245,200}

\renewcommand{\arraystretch}{1.45}
\small
\begin{tabularx}{\linewidth}{>{\bfseries}p{4.2cm} p{3.8cm} p{3.8cm} p{3.8cm} p{4.2cm}}
\toprule
  & \textbf{IB — Reg T}
  & \textbf{IB — Portfolio Margin}
  & \textbf{TR Model-Based}
  & \textbf{TR Simplified (this doc)} \\
\midrule
Core methodology
  & Flat regulatory rule (50\% of MV, Reg T / FINRA Rule 4210)
  & OCC scenario grid $\pm 15\%$ across 31 price points (TIMS); 15\% single-stock floor
  & Filtered historical simulation + EWMA volatility scaling; 99\% RVaR over 5-day horizon
  & Fixed $h_{\mathrm{IM}}$/$h_{\mathrm{MM}}$/$h_{\mathrm{ON}}$ per asset class; haircut $=$ LTV limit ($\mathrm{LTV} \leq 1-h$) \\
\rowcolor{gray!8}
Initial margin
  & \EUR{7{,}500} (50\% of MV)
  & \EUR{2{,}250} (15\% of MV)
  & \EUR{1{,}710} (11.4\% of MV)
  & $E_{\min} = h_{\mathrm{IM}} \times \mathrm{MV} = 30\% \times \EUR{15{,}000} = \EUR{4{,}500}$; LTV$_0 = 66.7\% \leq 70\%$ \\
$3{\times}$ leverage possible?
  & \textbf{No.} Requires \EUR{7{,}500} equity; max $2{\times}$ on \EUR{5{,}000}
  & \textbf{Yes.} \EUR{5{,}000} $>$ \EUR{2{,}250}
  & \textbf{Yes.} \EUR{5{,}000} $\gg$ \EUR{1{,}710}
  & \textbf{Yes.} LTV$_0 = 66.7\% \leq 70\% = 1 - h_{\mathrm{IM}}$ \\
\rowcolor{gray!8}
IM horizon
  & Until position is closed (50\% applies always)
  & Until position is closed (PM grid applies always)
  & 4 trading days, then MM regime
  & \textbf{1 trading day}, then MM regime \\
Maintenance / ongoing check
  & 25\% of current MV; Excess Liquidity $\geq 0$
  & 15\% scenario floor; Excess Liquidity $\geq 0$
  & LTV monitored at 75\% / 80\% / 90\% thresholds (graduated)
  & LTV $\leq 1 - h_{\mathrm{MM}} = 75\%$; single hard limit \\
\rowcolor{gray!8}
Overnight treatment
  & Overnight reverts to Reg T 50\% if intraday was lower
  & PM grid re-evaluated; same floor applies overnight
  & EOD haircut add-on raises LTV toward breach thresholds
  & LTV limit tightens to $1 - h_{\mathrm{MM}} - h_{\mathrm{ON}} = 70\%$ at EOD \\
LTV monitoring principle
  & Equity / MV $\geq 25\%$ (maintenance)
  & NLV $-$ PM\_req $\geq 0$
  & Graduated: monitoring 75\%, breach 80\%, liquidation 90\%
  & \textbf{Single limit:} $\mathrm{LTV} \leq 1 - h$; no graduated warnings \\
\rowcolor{gray!8}
Liquidation trigger
  & Excess Liquidity $< 0$ (binary)
  & Excess Liquidity $< 0$ (binary)
  & LTV $> 90\%$ (after warnings at 75\% and 80\%)
  & LTV $> 1 - h_{\mathrm{phase}}$ (immediate, no prior warning) \\
\rowcolor{highlight}
Intraday liquidation price $P^{*}$
  & $\EUR{5{,}000}/(100 \times 0.75) = \mathbf{\EUR{66.67}}$ \newline(on adjusted $2{\times}$ position)
  & $\EUR{10{,}000}/(150 \times 0.85) = \mathbf{\EUR{78.43}}$
  & $\EUR{10{,}000}/(150 \times 0.90 \times 0.90) = \mathbf{\EUR{82.30}}$
  & IM phase: $\EUR{10{,}000}/(150 \times 0.70) = \mathbf{\EUR{95.24}}$ \newline MM phase: $\EUR{10{,}000}/(150 \times 0.75) = \mathbf{\EUR{88.89}}$ \\
\rowcolor{highlight}
Price fall to intraday liquidation
  & $\mathbf{-33.3\%}$ (large buffer; $2{\times}$ position)
  & $\mathbf{-21.6\%}$
  & $\mathbf{-17.7\%}$
  & IM phase: $\mathbf{-4.8\%}$ \quad MM phase: $\mathbf{-11.1\%}$ \\
MTM frequency
  & Continuous (Excess Liquidity re-checked at each mark)
  & Continuous
  & Continuous (LTV re-computed per price update)
  & \textbf{Continuous} (LTV $= L / \mathrm{MV}_t$ at each price update) \\
\rowcolor{gray!8}
Asset-class differentiation
  & Partial (penny stocks, leveraged ETFs get higher rates; standard = flat 50\%)
  & Partial (single-stock 15\% vs.\ 10\% for broad indices)
  & Full (per-ISIN vol model; $h_{\mathrm{floor}}$ table override)
  & Full (per-class $h_{\mathrm{IM}}$, $h_{\mathrm{MM}}$, $h_{\mathrm{ON}}$ table) \\
Model complexity
  & None (rule-based)
  & Low (fixed scenario grid, OCC-licensed)
  & High (EWMA, filtered hist.\ simulation, correlation break, stress scenarios)
  & \textbf{Low} (table lookup; one LTV computation per price tick) \\
\bottomrule
\end{tabularx}

\bigskip
\small
\textbf{Notes.}
IB~Reg~T uses the adjusted $2{\times}$ position (100~shares, \EUR{10{,}000}~MV) because $3{\times}$ cannot be opened.
All other frameworks use 150~shares, \EUR{15{,}000}~MV, \EUR{10{,}000}~loan.
TR~Simplified uses $h_{\mathrm{IM}} = 30\%$, $h_{\mathrm{MM}} = 25\%$, $h_{\mathrm{ON}} = 5\%$ (other equity class).
TR~Model-Based liquidation price uses $h_{\mathrm{eff}} = 10\%$ and the 90\% LTV liquidation threshold.
Highlighted rows show the liquidation price and price fall — the primary practical differentiators.

\end{landscape}

\end{document}
